\documentclass[11pt,a4paper]{article}
\usepackage{fontspec}
\setmainfont{Noto Serif CJK SC}
\setsansfont{Noto Sans CJK SC}
\XeTeXlinebreaklocale "zh"
\XeTeXlinebreakskip = 0pt plus 1pt
\usepackage{amsmath,amssymb}
\usepackage[margin=2.1cm]{geometry}
\usepackage{graphicx,booktabs,array,caption,enumitem,titlesec,longtable}
\usepackage[section]{placeins}
\usepackage[hidelinks]{hyperref}
\sloppy
\linespread{1.25}
\setlength{\parindent}{2em}
\titleformat{\part}[display]{\sffamily\huge\bfseries\centering}{第\thepart 部分}{0.6em}{}
\titleformat{\section}{\sffamily\Large\bfseries}{\thesection}{0.8em}{}
\titleformat{\subsection}{\sffamily\large\bfseries}{\thesubsection}{0.8em}{}
\renewcommand{\figurename}{图}\renewcommand{\tablename}{表}\renewcommand{\contentsname}{目录}
\renewcommand{\thepart}{\arabic{part}}
\captionsetup{font=small,labelfont=bf,labelsep=quad}
\newcommand{\prop}[2]{\par\smallskip\noindent\textbf{#1}\quad #2\par\smallskip}
\newcommand{\jia}{[甲类]}\newcommand{\yi}{[乙类]}\newcommand{\bing}{[丙类]}
\newcommand{\fig}[3]{\begin{figure}[htbp]\centering\includegraphics[width=#2\textwidth]{fig/\detokenize{#1}}\caption{#3}\label{fig:#1}\end{figure}}
\setlist{nosep,leftmargin=2.2em}
\begin{document}
\begin{center}
{\sffamily\LARGE\bfseries 拉格朗日轨迹网络上的流动通道}\\[5pt]
{\sffamily\large 只凭示踪点位置的输运、掺混、动力学与热力学：理论、验证与应用（二维）}\\[4pt]
{\sffamily 完整报告}\\[8pt]
{\small 2026 年 10 月}
\end{center}

\noindent\textbf{摘要}\quad
本报告把此前以掺混为中心的工作扩展为一个统一框架，并给出全部检验。基础对象是一个拉格朗日轨迹网络（LTN）$G$：节点是轨迹，边是各帧互 kNN 对与 Delaunay 边的并集，每帧剥离 Delaunay 复形 $K(t)\subset G$ 是几何激活，成对伸长率是动力学激活。在 $K(t)$ 的三层余链上，只用示踪点位置与一个参考时刻的密度，给出质量、动量、能量、熵四类守恒或演化律的离散对应：账本给出连续方程；面元环量与加速度余链的旋度给出 Stokes 与 Kelvin 定理；$\rho a$ 余链的 Hodge 分解把压力（梯度通道）与粘性（旋度通道）分开，并由旋度通道功率反演粘度；三点映射与成对伸长率给出耗散；Dirac 算子给出线性声学并反演声速；$p/\rho^\gamma$ 给出熵与 RH 跳变；调和通道给出绕洞环量；$L_0$ 给出物质坐标下的扩散与掺混。

验证分三层。(1) 精确解：线性流、扩散 Lamb–Oseen、衰减 Taylor–Green、非均熵涡、线性声驻波、Sod 激波管、片层扩散、内重力波、Stokes 第一问题与 Lamb–Oseen 环形区；恒等式精确到舍入误差，等式在精确解上的误差为 0.05\%–2\%，收敛阶与噪声阈值均已测定。(2) DNS：激波–重气泡（无粘与粘性）、单模 RM（无粘、两档粘度）、声脉冲；压力与 DNS 相关不低于 0.994，粘度反演误差约 4\%，耗散的两种激活与 DNS 相差 1\%–2\%，声速 0.5\%，熵跳变与 RH 一致。(3) LHD 报告的九组算例重算并与 DTFE、SPH、FTLE、DBS、相遇数、LAVD 对比：账本密度比 DTFE 准 1–2 个数量级；"$\rho a$ 旋度份额"可以判断轨迹是否满足 Euler 方程；网络 LAVD 与精确 LAVD 的秩相关为 0.9997（双环流）与 0.81（涡对）。应用于有限雷诺数 RM：粘性削弱拉伸、导热削弱压缩，几何接管判据对 $Re$ 不敏感，掺混速率类判据在 $\mu=2\times10^{-3}$ 时推迟 0.25–0.35。正定二次型（涡拟能、耗散、熵产生）会把示踪点噪声计为物理量，报告给出由此得出的估计量设计准则、两个无量纲噪声数以及仍缺的验证。

\tableofcontents
\newpage

%=====================================================================
\part{理论}
%=====================================================================
\section{问题、范围与命题类别}
问题是：只给出一组示踪点的轨迹（加一个参考时刻的密度），能否不求速度场导数、不用网格，就描述可压缩含激波流动中的输运、掺混、压缩、旋转、压力、粘性耗散、声学与热力学，并让这些描述与动力学方程有可检验的等价关系。全文的计算只读取轨迹与参考密度；DNS 与精确解只用于打分。全部算例为二维。

命题分三类：\jia 是对任意位置成立的离散恒等式；\yi 是在写明的动力学或本构假设下成立的等式；\bing 是估计量或统计结论，精度只能由数值检验给出。第 \ref{sec:status} 节按这三类和参照层级逐项列出验证状态。

本构参数（粘度 $\mu$、扩散率 $\kappa$、导热系数 $k$、状态方程）不能从位置中直接读出。网络给出与这些参数相乘的几何二次型；反过来，也可以由某一通道的平衡关系反演参数。分子扩散与粘性不会让流体质点做随机运动，它们进入加速度（粘性力）与本构算子；对流在物质坐标中消失，体现为网络几何随时间的变化。

\section{拉格朗日轨迹网络与两种激活}
\prop{定义 2.1（LTN）}{$G=(V,E_G)$。节点是轨迹 $x_i(t)$，每个节点带一个在参考时刻固定的质量 $m_i=\rho_{0,i}A^*_i(t_0)/3$（$A^*_i$ 为星面积）。边集 $E_G=\bigcup_s\big(\mathrm{mkNN}_k(t_s)\cup\mathrm{Del}(K(t_s))\big)$；每条边是一对物质点，带历史 $r_e(t)=x_j(t)-x_i(t)$。}
\prop{定义 2.2（几何激活）}{每帧的外侧剥离 Delaunay 复形 $K(t_s)\subset G$（代码中逐边断言）。2-胞腔是活跃的物质三元组 $T$，历史为边矩阵 $D_T(t)=[x_j-x_i,\;x_k-x_i]$。}
\prop{定义 2.3（动力学激活）}{边信号为成对伸长率 $\varepsilon_e=\mathrm d\ln|r_e|/\mathrm dt$；AFRC 取窗口端点互 kNN 边上 $\varepsilon_e$ 的稳健中位数。}
\prop{定义 2.4（对称延拓）}{在对称面或壁面外放镜像轨迹，$G\to G\cup\sigma G$。对称流动中镜像轨迹本身是全域流动的轨迹；它消除对称线上凸包的钝角薄片三角形（不加时 Rayleigh 商偏高 32\%、吸收随机游走特征值偏高 15\%，加后 1.6\%）。对无滑移壁面，像轨迹取全平面问题的解（第 \ref{sec:exact} 节 Stokes 算例）。}
以粘性激波–重气泡为例：$|V|=31920$（含镜像），$|E_G|=355149$，$K(t)$ 平均 95519 条边，$E_G$ 中 49\% 的边从未在 $K(t)$ 中激活，只参与动力学激活。全部通道由同一个 LTN 类（继承 LHD 的 \texttt{trajnet.TrajectoryNetwork}）计算，与此前各分散脚本的结果在 0.2\% 以内一致（图 \ref{fig:fig12.jpg}a）。

\section{余链、算子与度量}
在 $K(t)$ 上取节点–边、边–面关联矩阵 $B_1$、$B_2$（$B_1B_2=0$），Hodge 星 $\star_0=\mathrm{diag}(A^*_i/3)$（或质量 $m_i$），$\star_1=\mathrm{diag}(w_e)$，$w_e=\tfrac12(\cot\alpha_e+\cot\beta_e)$，$\star_2=\mathrm{diag}(1/A_T)$。边余链由节点属性构造：
\begin{equation}
u_e=\tfrac12(u_i+u_j)\cdot r_e,\quad a_e=\tfrac12(a_i+a_j)\cdot r_e,\quad f_e=\tfrac12(\rho_ia_i+\rho_ja_j)\cdot r_e,
\end{equation}
$u_i$、$a_i$ 为位置的中心一阶、二阶差分，$\rho_i=\rho_{0,i}/J_i$ 来自物质账本。
\prop{命题 3.1（Hodge 分解）\jia}{$C^1=\mathrm{im}\,B_1^{\top}\oplus\mathrm{im}\,\star_1^{-1}B_2\oplus\mathcal H$，在 $\star_1$ 度量下正交；梯度势 $\phi$ 在节点上，旋度势 $\chi$ 在面上；$\dim\mathcal H=\beta_1$。}
\prop{命题 3.2（三点映射）\jia}{$F_T=D_T(t+\Delta)D_T(t-\Delta)^{-1}$，$\nabla u_T\approx\log F_T/(2\Delta)$（$2\times2$ 精确矩阵对数），$\mathrm{tr}\log F_T=\ln[A_T(t+\Delta)/A_T(t-\Delta)]$ 精确成立。}
\prop{命题 3.3（物质账本）\jia}{$\ln J_i$ 由固定三角形上的星面积比逐步链接，激波块内推迟重锚（C 通道）。}
\prop{定义 3.4（Dirac 算子）}{在 $C^0\oplus C^1$ 上 $D=\begin{pmatrix}0&\delta_1\\ d_0&0\end{pmatrix}$，$d_0=B_1^{\top}$，$\delta_1=-\star_0^{-1}B_1\star_1$，$D^2=\mathrm{diag}(L_0,L_1^{\downarrow})$。}
\noindent\textbf{物质形式原则。} 凡能用物质形式（面积率、三点映射、环量）表示的量，都不用速度差分：位置差分得到的速度带 $O(u'''\Delta t^2)$ 误差，除以 $h$ 被放大。速度余链的 Galerkin 散度与 DNS 的 Spearman 相关只有 0.37–0.76，物质面积率为 0.54、0.94、0.91。

\section{通道总览：物理意义、构造理由与对标}\label{sec:overview}
各通道不是并列的指标，而是同一离散外微分结构的不同层：点层放质量、熵、标量；边层放速度、加速度、$\rho a$；面层放环量；$B_1$、$B_2$ 连接相邻层。每条守恒律或演化律恰好落在一层或一个 Hodge 通道上。

\textbf{质量与压缩（点层）。} 物质微团面积变化即密度变化；描述激波、压缩波、膨胀扇、气泡压扁与导热膨胀；固定三角形面积比是行列式，精确成立；不可压时恒为 1，激波在一帧内完成 $\eta$ 倍压缩，$\eta$ 直接进入 RH 关系。对标：PIV 散度、DTFE、SPH 密度、Cauchy–Green 行列式、Ducros 型激波传感器。

\textbf{旋转与 Kelvin（面层）。} 面元环量即涡量积分，加速度环量即环量变化率；描述涡片、卷起、涡核、剪切层；机理为斜压、粘性扩散与 Crocco 定理；Stokes 与 Kelvin 在物质三角形上离散精确；激波扫过界面时一次性沉积环量（Richtmyer 冲量模型、Samtaney–Zabusky 模型）。对标：PIV 环量法、Q/$\lambda_2$、LAVD。

\textbf{梯度通道一：压力。} $\rho a$ 无粘时严格是梯度（$a$ 本身含斜压旋度，不是梯度），Hodge 投影等价于余切权 Poisson 最小二乘，残差即诊断量；激波冲量的空间积分等于 RH 动量关系给出的 $\Delta p$。对标：PIV/PTV 压力重建（全向积分、Shake-The-Box、可压缩 PIV 的等熵假设）。

\textbf{旋度通道：粘性。} 二维常粘度 $\nabla\cdot\tau=\tfrac43\mu\nabla\theta-\mu\,\mathrm{curl}\,\omega$，旋度部分只含粘性；其对速度的功为 $-\mu\int\omega^2$，只需一阶导数即可反演粘度；激波内粘性是数值的。对标：PDE-FIND、PINN、动态涡粘、格式耗散评估。

\textbf{应变与耗散（两种激活）。} 三点映射给出完整 $\nabla u$；成对伸长率因转动不改变长度天然只测 $S$，多边拟合抗噪，且正是 AFRC 的边信号；激波处 $\int\theta^2$ 不收敛，耗散改由熵跳变表示。对标：tetrad 方法、FTLE/Cauchy–Green、动态极分解、成对扩散统计。

\textbf{声学（Dirac）。} 点信号（密度）与边信号梯度部分相互驱动；$D^2$ 给出波动方程，反对称结构保证能量守恒，可反演声速；激波由账本与 RH 接管。对标：Kovasznay 分解、可压缩湍流 Helmholtz 分解、拓扑 Dirac 信号处理。

\textbf{熵与第二定律（点层）。} 压力（动量通道）与账本密度闭合热力学；无粘时熵是物质标签，只有激波改变它；导热熵产生是 $\ln T$ 的拉普拉斯二次型，第二定律即其正定性。对标：等熵假设推温度、DNS 熵收支、熵型激波传感器。

\textbf{调和通道（拓扑）。} 绕洞环量，$\beta_1$ 为洞数；周期是拓扑不变量，能量份额随度量变化。对标：环量测量、编织与拓扑熵、持续同调、网络调和投影。

\textbf{扩散与掺混（点层 $L_0$）。} 物质坐标下对流消失，变形进入电导，Pe 即 Markov 时间；分接触、交换、状态三层。对标：片层理论、离析强度、mix-norm、传递算子、EMST 掺混模型、RM 掺混度量。

\textbf{结构与激活统计。} AFRC/Ollivier 曲率给出屏障与相干集，Forman 曲率是 $L_1$ 的 Bochner–Weitzenböck 余项；边出生率区分仿射翻转与搅拌新接触。对标：FTLE/LCS、LAVD、DBS、动态 Laplacian、$\varepsilon$-网络、相遇数。

\textbf{可压缩与激波的总差别。} (1) 不变测度是质量，面积加权会把压缩误读为掺混、膨胀或耗散；(2) 速度梯度通道、声学通道、熵通道只在可压缩流中非空，并经斜压项与涡通道耦合；(3) 激波在示踪点尺度上是间断，正定型不收敛，激波强度与不可逆性由账本跳变加 RH 表示，动量通道仍可用；(4) 激波造成大量 Delaunay 翻转，组合拓扑量必须过压缩滤网或按参考距离过滤。

%=====================================================================
\part{各通道的命题与结果}
%=====================================================================
\section{质量与压缩}
\prop{命题 5.1 \jia}{账本在物质星上给出 $\mathrm D\ln\rho/\mathrm Dt=-\theta$，$\theta_i=\mathrm d\ln A^*_i/\mathrm dt$。}
RM 界面带压缩：单次激波 0.75、再激波 1.16，与 LHD 报告的 0.69–0.75、1.11–1.19 一致；DNS 密度在界面给出 0.84，是接触面抹平所致。精确解与 DTFE 对比见第 \ref{sec:exact}、\ref{sec:lhd} 节。

\section{旋转与 Kelvin 定理}
\prop{命题 6.1 \jia{} 与 \yi}{Stokes：$\Gamma_T=(B_2^{\top}u)_T$。Kelvin：对物质三元组 $(B_2^{\top}a)_T=\mathrm d\Gamma_T/\mathrm dt$；无粘时等于 $\int_T(\nabla\rho\times\nabla p)/\rho^2\,\mathrm dA$。}
激波–重气泡（无粘）：面元涡量与 DNS 相关 0.98–1.00；$\int\omega^2$ 为 1.672/1.702、6.004/6.072、4.794/4.841（网络/DNS，$t=0.3,0.6,1.0$）；整团环量变化率 $-0.655/-0.640$、$-0.545/-0.604$、$-0.080/-0.059$，分箱相关 0.96、0.91、0.72（图 \ref{fig:fig6.jpg}、\ref{fig:fig7.jpg}）。

\fig{fig6.jpg}{1.0}{激波–重气泡（无粘）：由位置重建的压力、加速度余链的旋度（Kelvin）、速度余链的流函数与物质面积率；DNS 场仅作参照。}
\fig{fig7.jpg}{1.0}{运动学–动力学对应的定量检验：压力误差、动能的 Hodge 分解、加速度与 $\rho a$ 的旋度份额、Kelvin 积分、耗散恒等式两项、散度的两种算法。}

\section{动量：压力与粘性}
\prop{命题 7.1 \yi}{$\rho a$ 余链的梯度部分对应 $-\nabla p+\tfrac43\mu\nabla\theta$，余正合部分对应 $-\mu\,\mathrm{curl}\,\omega$；$p=-\phi_f+\tfrac43\mu\theta+\mathrm{const}$。}
\prop{命题 7.2 \yi}{$\langle f_c,u\rangle_{\star_1}=-\mu\int\omega^2+$ 边界项；$\mu_E=-\langle f_c,u_c\rangle/\int\omega^2$；逐面元 $\chi_f=-\mu\omega_T$。}
压力（无粘激波–重气泡，19 个时刻）：相关不低于 0.994，RMS/std 为 2\%–16\%（全部）、0.5\%–9\%（去激波块）；需每帧一个参考常数。无粘时 $\rho a$ 旋度份额在 $t\le0.8$ 不超过 0.3\%，$t=1.38$ 为 4.4\%。粘度（图 \ref{fig:fig8.jpg}）：粘性激波–重气泡（真值 $5\times10^{-4}$）后期 5.1–5.2$\times10^{-4}$，激波作用期 5.8–8.6$\times10^{-4}$；无粘算例 1.1–2.9$\times10^{-4}$（数值有效粘度）；逐面元 $\chi_f=-\mu\omega$ 斜率 $5.85\times10^{-4}$、相关 $-0.91$；旋度通道功率与 DNS 物理粘性力相差约 2\%。力矩形式 $\mathrm{curl}(\rho a)=\mu\nabla^2\omega$ 失败（需三阶导数）。

\fig{fig8.jpg}{1.0}{粘性通道（激波–重气泡）：由位置反演粘度、动能收支的 Hodge 通道、耗散、粘性力的空间结构、旋度通道面势 $\chi=-\mu\omega$。}

\section{能量：耗散的两种激活}
\prop{命题 8.1 \jia}{$2S{:}S-\tfrac23\theta^2=\omega^2+\tfrac43\theta^2-4\det\nabla u$；$\int\det\nabla u$ 只贡献边界积分，对分片线性映射离散精确。}
\prop{命题 8.2 \yi}{几何激活：$\Phi_T=\mu[2S{:}S-\tfrac23\theta^2]_T$；动力学激活：$\varepsilon_e=\hat r_e^{\top}S\hat r_e$ 逐节点最小二乘；动能收支 $\langle f,u\rangle=\langle f_g,u_g\rangle+\langle f_c,u_c\rangle$（交叉项约 $10^{-8}$）。}
几何激活与 DNS 光滑区 $\int\Phi$ 之比：$t=0.3$–0.62 为 0.97–1.00，$t\ge1.1$ 为 0.98–1.004，$t\approx0.7$–1.0 为 0.85–0.93；不计边界项偏高约 16\%。两种激活的应变逐点相关 0.998–0.999，动力学激活与 DNS 相差 1\%–2\%（图 \ref{fig:fig12.jpg}b）。AFRC 所用边信号的二阶统计即粘性耗散。

\fig{fig12.jpg}{1.0}{同一个 LTN 上的回归与两种激活：(a) 统一实现与此前脚本之比；(b) 耗散的两条途径；(c) $G$ 的激活统计。}

\section{声学：Dirac 通道}
\prop{命题 9.1 \yi}{线性等熵时 $a=-c^2\nabla\ln\rho$，$c^2=-\langle a,a\rangle/\langle a,d_0\ln\rho\rangle$（账本梯度含噪，置于因变量一侧）；$\partial_t(\rho',u)=-(\rho_b\delta_1u,\;c^2\rho_b^{-1}d_0\rho')$ 为 Dirac 型方程，$D^2$ 给出 $\omega^2=c^2k^2$。}
声脉冲 DNS（图 \ref{fig:fig9.jpg}）：$c^2=1.005$（反向回归会得到 0.61–0.75 的衰减偏差）；速度旋度份额约 $7\times10^{-7}$；Dirac 前推与 DNS 相差 0.4\%、1.1\%、6.4\%（$t=0.1,0.2,0.3$），能量守恒到 $10^{-13}$。

\fig{fig9.jpg}{1.0}{(a) 三点映射的耗散闭合（粘性激波–重气泡）；(b)–(f) 声脉冲：声速反演、径向剖面、声学能量与 Dirac 前推。}

\section{熵与第二定律}
\prop{命题 10.1 \yi}{$s=c_v\ln(p/\rho^\gamma)$；无粘光滑区 $\mathrm Ds/\mathrm Dt=0$；激波处 $\Delta s(\eta)$ 由 RH 给出，$\eta$ 只用账本。}
\prop{命题 10.2 \yi}{$\rho T\,\mathrm Ds/\mathrm Dt=\Phi+\nabla\cdot(k\nabla T)$；熵产生 $\int\Phi/T+k\sum_ew_e(\Delta\ln T)_e^2\ge0$。}
熵场与 DNS 相关 0.99、0.98；重/轻气体熵差 $-3.75$、$-3.76$（DNS $-3.77$、$-3.75$，初值 $-3.85$）；光滑区 0.4 时间单位内 $\Delta s$ 的 RMS 为 0.045–0.086（DNS 0.076–0.127）；激波熵跳变：重气体 0.150/0.151/0.148（RH/网络/DNS），轻气体 0.063/0.049/0.066。粘性算例温度噪声 0.8\%–1.5\% 使导热熵产生被放大 2–5 倍；低通 $\ell=2h$ 后重气体熵变化率与 DNS 相差 10\%–20\%（图 \ref{fig:fig10.jpg}）。

\fig{fig10.jpg}{1.0}{熵通道：熵场、激波熵跳变的三条途径、光滑区守恒、粘性算例的熵收支与第二定律。}

\section{调和通道}
\prop{命题 11.1 \yi（连续极限）}{环形区内径向对称涡量时 $h=C/(2\pi r)\,e_\theta$，$C=\int\Gamma(r)r^{-1}\mathrm dr/\ln(r_{out}/r_{in})$；加速度调和周期为 $\mathrm dC/\mathrm dt$（非单连通区域的 Kelvin 定理）。}
扩散 Lamb–Oseen 环形区：只剥离与洞相邻的长三角形后 $\beta_1=1$；速度调和周期与理论相差不超过 0.8\%，加速度 3\%–5\%；能量份额随度量变化（95\% 与 78\%），应报告周期（图 \ref{fig:fig11.jpg}）。

\fig{fig11.jpg}{1.0}{(a)(b) 正定型的噪声偏置与四种处理；(c)–(f) 扩散 Lamb–Oseen 环形区的调和通道。}

\section{扩散与掺混}
\prop{命题 12.1 \jia}{以质量为点权、余切电导为边权，网络扩散是物质坐标扩散方程 $\rho_0\partial_tc=\nabla_X\cdot(kJC^{-1}\nabla_Xc)$ 的 P1 离散。}
\textbf{压缩滤网（Sod 精确运动学）。} 组合量失败：$\varepsilon$-网络平均度 +126\%，聚类系数 +24\%，Delaunay 边更替 24\%，无权 Dirichlet 能量 $-32\%$ 至 $+72\%$；余切电导加质量的量通过：Dirichlet 能量不超过 1\%，余切同配 $-1.2\%$，远程接触为 0。

\textbf{激波–重气泡。} Pe 扫描给出 $t=1.4$ 的掺混增强 1.44、1.38、1.31、1.23（$\Gamma=3\times10^{-5}$ 至 $10^{-3}$）；网络掺混时钟预测高 Pe 极限 1.34–1.42；逐点解释"流动在哪里加速了掺混"：网络时钟 Spearman 0.83/0.59，AFRC 0.46/0.26，远程接触 0.08/0.34，距离约 0（图 \ref{fig:fig1.jpg}、\ref{fig:fig2.jpg}）。

\textbf{RM。} 重写的 DNS 复现掺混文档的参照接管时刻（几何 3.39 对 3.43，时钟 2.25 对 2.19，再激波 8.43 对 8.45、5.89 对 5.9）。尖锐标签割长分箱的 $\langle\lambda^2\rangle$ 与标记曲线相差不超过 0.1（单次激波，到 $t=6$）。以 $\lambda^2/J$ 的算术平均为判据接管最早（单次 1.96/1.80，再激波 4.10/3.78，标记曲线/网络）。混合质量判据随 $\Gamma$：单次 3.84、4.00、5.18；再激波 9.05、10.2、大于 12；$\Gamma\to0$ 时趋向片层极限 3.6、8.6（图 \ref{fig:fig3.jpg}–\ref{fig:fig5.jpg}）。

\fig{fig1.jpg}{1.0}{激波–重气泡：网络掺混场（物质标签、网络掺混时钟、局部混合度、流动带来的掺混增量、远程接触、AFRC）。}
\fig{fig2.jpg}{1.0}{掺混通道的定量检验：压缩滤网、接触与交换势、Pe 扫描、掺混增强、吸收随机游走、逐点解释力。}
\fig{fig3.jpg}{0.95}{单次激波 RM 的网络场：标签与标记曲线、账本压缩、网络拉伸、网络掺混时钟、混合度、掺混增量。}
\fig{fig4.jpg}{0.95}{再激波 RM 的网络场。}
\fig{fig5.jpg}{1.0}{RM 接管判据：拉伸与压缩因子、界面长度、各判据的接管时刻、混合质量判据随 Pe 的变化。}

\textbf{有限雷诺数 RM（图 \ref{fig:fig13.jpg}，表 \ref{tab:re}）。} 粘度反演：$\mu=5\times10^{-4}$ 时 5.1–5.4$\times10^{-4}$（$t\le3$），$\mu=2\times10^{-3}$ 时 1.88–2.04$\times10^{-3}$（$t\le1.5$），之后下降（边界项，见第 \ref{sec:exact} 节 Stokes 算例）。到 $t=6$，$\ln L/L_0$ 为 1.50、1.08、0.87，$\langle-\ln J\rangle$ 为 0.70、0.47、0.37：粘性削弱拉伸，导热削弱压缩（冷重气体被加热膨胀；双组分 RM 中对应效应来自互扩散）。
\begin{table}[h]\centering\caption{接管时刻随粘度（标记曲线参照；最后一行为 LTN 反事实扩散，$\Gamma=10^{-4}$）}\label{tab:re}
\begin{tabular}{lccc}\toprule
判据 & 无粘 & $\mu=5\times10^{-4}$ & $\mu=2\times10^{-3}$\\\midrule
几何：$\ln L$ 超过压缩 & 3.39 & 3.33 & 3.38\\
时钟 $\alpha=1$ & 2.25 & 2.25 & 2.50\\
速率 $\alpha=1$ & 1.96 & 1.95 & 2.21\\
速率 $\alpha=0$ & 3.01 & 3.05 & 3.35\\
混合质量 $\alpha=1$ & 4.00 & 4.13 & 4.32\\\bottomrule
\end{tabular}\end{table}
\fig{fig13.jpg}{1.0}{有限雷诺数 RM：粘度反演、两种激活的耗散、接管时刻、拉伸与压缩。}

\section{激活统计}
粘性激波–重气泡每节点每单位时间新激活 1.6–3.3 条边，其中参考距离大于 $3h$ 的新接触比例由 $t=0.3$ 的 3\% 升到 0.6 的 17\%、$t\ge1.0$ 的 43\%–54\%；RM 由 14\% 升到 41\%–62\%。早期是压缩造成的仿射翻转，后期是搅拌。

%=====================================================================
\part{验证}
%=====================================================================
\section{参照的层级}
\begin{itemize}
\item \textbf{R1 精确解}（与数值格式无关）：线性流、Lamb–Oseen、Taylor–Green、非均熵涡、声驻波、Sod、片层、内重力波、Stokes 第一问题、Lamb–Oseen 环形区。
\item \textbf{R2 同一 DNS 的场量}：示踪点由 DNS 速度插值推进，只能检验网络算子与离散流场是否自洽，分不开方法误差与 DNS 数值误差。
\item \textbf{R3 同一 DNS 中推进的标记曲线}：与 R2 同源。
\item \textbf{R4 内部一致性}：两种激活一致、回归、正交性、离散能量守恒、Galilean 不变。
\end{itemize}

\section{精确解验证}\label{sec:exact}
全部计算经同一个 LTN 类，不用任何 DNS；轨迹由解析解给出或对解析速度场做 RK4 积分（细时间网格必须与帧时刻对齐，否则时间错位会被二阶差分放大为虚假加速度）。验证矩阵见图 \ref{fig:fig15.jpg}，场图见图 \ref{fig:fig16.jpg}、\ref{fig:fig17.jpg}。
\begin{itemize}
\item \textbf{线性流}（转动、纯应变、简单剪切、各向同性膨胀、一般）：账本膨胀率、三点映射 $\nabla u$、耗散为 $10^{-13}$–$10^{-12}$（舍入误差）；涡量、成对应变、Kelvin 力矩、压力为 $10^{-6}$–$10^{-5}$（位置差分的 $O(\Delta t^2)$）。
\item \textbf{扩散 Lamb–Oseen}（$h=0.005\to0.0025$）：$\mu_E$ $+0.26\%\to-0.04\%$，压力 $0.11\%\to0.05\%$，涡量 $1.4\%\to0.7\%$，耗散（成对）$-0.53\%\to-0.21\%$。粘性力矩（加速度旋度）误差 110\%$\to$55\%：弱力矩叠加强向心加速度时，P1 边积分的 $O(h^2\nabla^2a)$ 误差超过力矩本身；减小 $\Delta t$ 无效。
\item \textbf{Taylor–Green}：$\mu_E$ $-0.32\%$，压力 0.22\%，涡量 1.1\%，耗散（成对）$-0.40\%$、（三点映射）$-4.9\%$。
\item \textbf{非均熵涡}（变密度定常 Euler）：压力 0.10\%，熵 0.10\%，账本 $1.4\times10^{-4}$；$\rho a$ 旋度份额 $5\times10^{-7}$。
\item \textbf{线性声驻波}：$c^2$ $-0.13\%$，旋度份额 $1.5\times10^{-6}$，Dirac 前推 0.08\%/0.17\%（$t=0.05/0.1$），能量漂移 $4\times10^{-10}$；账本相对波幅 0.2\%–2.5\%，随时间累积。
\item \textbf{Sod}：跨激波 $\eta$ 与 RH 熵跳变精确到机器精度；压力 1.9\%–2.1\%（含激波）、0.1\%–0.8\%（去激波块）；密度 1.2\%（误差在接触面）；稀疏波膨胀率 8\%–11\%（折点被星平均抹平）。
\item \textbf{片层}：纯应变剖面 2.4\%、混合质量 2.1\%；一维压缩 $\alpha=1$、$\alpha=0$ 剖面 1.5\%、1.6\%。面积形式在均匀压缩中恰好等于 $\kappa$ 为常数（$\alpha=0$）的解，用于 $\alpha=1$ 时偏差 19\%；错误在本构不一致。
\item \textbf{内重力波}（强斜压力矩）：Kelvin 力矩与 $-N^2\partial\zeta/\partial x$ 相差 1.4\%，由位置反演的 $N^2$ 误差 0.08\%。
\item \textbf{Stokes 第一问题}（壁面）：壁面涡量 0.26\%/0.28\%，耗散 0.05\%；$\mu_E$ 在墙作自由边界时 $-120\%$，墙外放像轨迹后 $+6\%$。旋度通道恒等式在边界处多出 $\mu\times$ 边界涡量的项。
\item \textbf{收敛阶}（Taylor–Green）：涡量 1，压力约 2，成对耗散约 1.9，三点映射耗散约 1（图 \ref{fig:fig18.jpg}a）。
\item \textbf{噪声阈值}（误差 5\%）：加速度类 $\sigma/h\approx3\times10^{-4}$（$N_a\lesssim0.05$），涡量 $5\times10^{-4}$，三点映射耗散 $3\times10^{-3}$，成对耗散 $5\times10^{-3}$（$N_g\lesssim0.2$）（图 \ref{fig:fig18.jpg}b）。
\end{itemize}
\fig{fig15.jpg}{1.0}{精确解验证矩阵（每格为相对精确解的误差）、Sod 剖面与 Lamb–Oseen 加密收敛。}
\fig{fig16.jpg}{1.0}{精确解场图 I：Lamb–Oseen、Taylor–Green、非均熵涡的涡量、压力、熵与耗散（精确 / 网络 / 误差）。}
\fig{fig17.jpg}{1.0}{精确解场图 II：声驻波、片层掺混、内重力波斜压力矩、Stokes 壁面涡量。}
\fig{fig18.jpg}{0.95}{Taylor–Green 的收敛阶与位置噪声阈值。}

\noindent 误差形态：涡量与耗散的误差是示踪点尺度的随机斑点，加密即收敛；压力误差是光滑系统量，集中于涡核，是 Poisson 重建的二阶截断误差，不能靠平滑去除；熵误差继承自压力。

\section{DNS 上的敏感性}
粘性激波–重气泡，$t=1.2$（图 \ref{fig:fig14.jpg}）：
\begin{table}[h]\centering\small\caption{敏感性（$\mu_E$ 单位 $10^{-4}$，真值 5.0）}
\begin{tabular}{lcccc}\toprule
变动 & $\mu_E$ & 三点映射 $\Phi$/DNS & 成对 $\Phi$/DNS & 压力误差\\\midrule
基准 & 5.18 & 1.00 & 0.99 & 0.8\%\\
帧间隔加倍 & 5.11 & 1.01 & 0.98 & 2.6\%\\
间距 $\sqrt2h$ & 5.44 & 1.01 & 0.98 & 1.0\%\\
噪声 $0.2\%h$ & 5.39 & 1.02 & 0.99 & 1.2\%\\
噪声 $1\%h$ & 5.57 & 1.34 & 1.02 & 4.0\%\\
噪声 $4\%h$ & 2.67 & 6.5 & 1.36 & 16\%\\\bottomrule
\end{tabular}\end{table}
\fig{fig14.jpg}{1.0}{DNS 上的帧间隔、示踪点间距与位置噪声敏感性。}

\section{噪声原理与估计量设计准则}
\prop{命题 17.1 \bing}{正定型 $Q(x)=x^{\top}Lx$ 满足 $\mathbb E\,Q(x+n)=Q(x)+\mathbb E\,n^{\top}Ln\ge Q(x)$：独立噪声对涡拟能、耗散、熵产生给出正偏差。}
\begin{enumerate}
\item 用物质形式，不用速度差分。
\item 用能量或积分形式，不用高阶导数（$\mu_E$ 成功，力矩形式失败）。
\item 回归时把含噪一侧放在因变量（$c^2$）。
\item 正定型须去偏：分样交叉估计消去独立噪声（导热熵产生超出量 80\%–380\% $\to$ 30\%–65\%）；共模误差（同一插值速度场、同一参考常数）无法由内部统计去除；GCV 失败（$\ell^*=0.002$–0.003）；固定 $\ell=2h$ 偏低 20\%–28\%。
\item 多边拟合抗噪：成对伸长率在 $\sigma/h=3\times10^{-3}$ 时耗散误差仍为 0.7\%。
\item 通道可用性由 $N_g=\sigma/(|\nabla u|h\Delta t)$ 与 $N_a=\sigma/(|a|\Delta t^2)$ 决定。
\end{enumerate}

%=====================================================================
\part{与 LHD 报告算例和主流方法的对比}\label{sec:lhd}
%=====================================================================
\section{重算设置}
LHD 报告的八个算例（双环流、Lamb–Oseen 涡对、Sod、柱面爆炸波、激波–等熵涡、激波–重气泡、粘性激波–重气泡、单模 RM）加粘性 RM 共九组，全部经同一个 LTN 类重算。对比方法：报告表 7 的账本、SPH、DTFE 密度误差与 T3、T7 数值；同一组示踪点上新算的 DTFE、FTLE、扩散屏障强度 DBS（Haller–Karrasch–Kogelbauer，最小二乘 $F$）、相遇数（Rypina–Pratt）；LTN 自身的 AFRC 屏障强度与由面元涡量积分的 LAVD。这里的 RM 设置（$a_0=0.1$、随界面参考系、界面带示踪点）与报告不同，RM 数值不能直接对照。

\section{质量通道}
\begin{table}[h]\centering\small\caption{密度相对误差（中位 / 95 分位；爆炸波为 $J$）}
\begin{tabular}{lccc}\toprule
算例 & LTN 账本（本次） & DTFE（本次） & 报告：账本 / SPH / DTFE（中位）\\\midrule
双环流 & 1.2\% / 5.8\% & 32\% / 256\% & 1.3\% / 11.4\% / 32.1\%\\
Lamb–Oseen 涡对 & 0.53\% / 16\% & 26\% / 151\% & 0.41\% / 6.4\% / 26.3\%\\
Sod & $\approx0$ / 0.81\% & 0.25\% / 39\% & $<10^{-6}$ / 1.9\% / 0.20\%\\
爆炸波 & 0.34\% / 2.7\% & 12\% / 71\% & 95 分位：v4 账本 2.1\%\\
激波–等熵涡 & 0.36\% / 4.7\% & 18\% / 109\% & 0.34\% / 3.4\% / 16.8\%\\
激波–重气泡 & 0.57\% / 21\% & 20\% / 99\% & 0.67\% / 7.1\% / 25.9\%\\\bottomrule
\end{tabular}\end{table}
LTN 版账本复现了报告数值，在每个算例上比 DTFE 准 1–2 个数量级、比报告中的 SPH 准约一个数量级。RM 中位误差 2.4\%–3.8\%（报告 0.73\%）来自设置不同与 DNS 接触面的数值扩散。

\section{动力学通道}
\begin{table}[h]\centering\small\caption{新增动力学通道（末时刻附近）}
\begin{tabular}{lccccc}\toprule
算例 & 涡量相关 & 压力误差 & 熵相关 & $\rho a$ 旋度份额 & $\mu_E$\\\midrule
双环流 & 1.000 & — & — & \textbf{0.20} & —\\
L–O 涡对 & — & — & — & 0.02 & —\\
Sod & 无旋 & 0.13\%–0.35\% & 0.994 & 0.002–0.01 & $\approx0$\\
激波–等熵涡 & 0.999 & 1.5\%–1.8\% & 0.11–0.31 & 0.0005–0.001 & $\approx0$\\
激波–重气泡 & 0.995–0.996 & 0.9\%–9\% & 0.98–0.99 & 0.001–0.015 & 1.7–2.6$\times10^{-4}$\\
粘性激波–重气泡 & 0.999 & 0.5\%–7\% & 0.997–0.998 & $<0.001$ & 5.9–7.0$\times10^{-4}$\\
RM & — & 4.5\%–5\% & 0.95–0.98 & 0.005–0.024 & $\approx0$\\
RM $\mu=5\times10^{-4}$ & — & 3.8\% & 0.99 & 0.006 & 4.1$\times10^{-4}$\\\bottomrule
\end{tabular}\end{table}
三点发现：(1) "$\rho a$ 旋度份额"可判断轨迹是否满足 Euler 方程：双环流是运动学模型而非 Euler 解，份额 0.20，真正的 Euler 流动为 0.001–0.02。(2) 激波–等熵涡的熵相关低，因 $M_s=1.3$ 弱激波熵跳变约 $10^{-2}$，被压力误差淹没。(3) 爆炸波从 $t=1$ 开始而没有给出该时刻密度，用 $\rho_0=1$ 代入时旋度份额 0.43，说明参考密度设错会被该诊断检出。

\section{输运结构}
AFRC 与 FTLE 的秩相关为 0.84–0.97（Sod 为 0.64）。DBS 在强拉伸与折叠区出现病态点（初末帧之间的最小二乘 $F$ 失真，个别点 $J<0$），与其他方法相关只有 0.3–0.9；长时间窗的 DBS 需要逐帧链接 $F$ 或改用三点映射。Sod 中 FTLE 与 DBS 几乎不相关（0.09），因为 DBS 乘了 $J$、把压缩计入。网络 LAVD 与精确 LAVD：双环流秩相关 0.9997、前 10\% Jaccard 0.98；涡对 0.81、0.96（报告 T7 以动态 Laplacian 识别涡核的 Jaccard 为 0.82，目标不同，仅作参照）。
\fig{fig19.jpg}{1.0}{LHD 报告九个算例的 LTN 通道场图（账本与 DTFE 密度、面元涡量、压力、熵或网络 LAVD）。}
\fig{fig20.jpg}{1.0}{输运结构：FTLE、DBS、相遇数与 LTN 的 AFRC、LAVD（同一组示踪点、同一时间窗，画在初始位置）。}
\fig{fig21.jpg}{1.0}{定量对比：密度误差（含报告表 7）、动力学通道、输运结构方法之间的秩相关。}

\section{未重算的报告条目}
T1 激波检出（LTN 使用同一检测）、T4 同边界环量线积分与 Samtaney–Zabusky 比较、T6 质量输运、T7 动态 Laplacian、T8 扩散与 SPH 比较（同一 P1 算子，已回归）、T9 RM 振幅与 Richtmyer 模型（设置不同）、T10 噪声（由图 \ref{fig:fig14.jpg}、\ref{fig:fig18.jpg} 覆盖）。成对伸长率在爆炸波中与三点映射的应变相关仅 0.49–0.62，原因是按固定半径挑边；应改为按 $K(t)$ 邻接挑边。

%=====================================================================
\part{状态与展望}
%=====================================================================
\section{逐项状态}\label{sec:status}
\begin{longtable}{p{2.8cm}p{1.5cm}p{5.2cm}p{5.6cm}}
\caption{各结论的参照与状态}\\\toprule
结论 & 参照 & 状态 & 主要条件与风险\\\midrule\endhead
质量（账本） & R1 R2 & 线性流精确；Sod 跨激波精确；比 DTFE 准 1–2 量级 & 接触面附近几个示踪点宽的抹平\\
涡量 / Kelvin & R1 R2 & 涡量 0.7\%–1.4\%（一阶收敛）；强力矩 1.4\% & 弱力矩叠加强向心场时不可用\\
压力 & R1 R2 & 精确解 0.05\%–0.2\%（二阶）；DNS 相关 $\ge0.994$ & 参考常数；$N_a\lesssim0.05$\\
粘度 $\mu_E$ & R1 R2 & 精确解 0.04\%–0.3\%；DNS 约 4\% & 网络边界须无涡量（壁面放像轨迹）\\
耗散 & R1 R2 R4 & 成对 0.2\%–0.5\%；两种激活一致 & 三点映射对噪声与曲率敏感\\
声速 / Dirac & R1 & $c^2$ 0.13\%；前推 0.1\%–0.2\% & 线性范围\\
熵 / RH & R1 R2 & Sod 精确；非均熵涡 0.1\% & 弱激波不可分辨\\
第二定律 & R2 & 上下界 & 共模误差\\
调和 & R1 & 周期 0.8\%、3\%–5\% & 仅解析流动\\
掺混 & R1 R3 & 片层 1.5\%–2.4\%；复现文档参照 & 质量度量与本构一致；信息极限约 $2h$\\
结构 & R4 & AFRC 与 FTLE 0.84–0.97；LAVD 0.81–0.9997 & 与 Hodge 通道的定量关系未写出\\
有限 $Re$ RM & R3 R4 & 两档 $Re$ & 构型特殊（冷重气体）\\\bottomrule
\end{longtable}

\section{仍缺的验证与计划}
\begin{enumerate}
\item 独立 DNS：换数值格式或加密网格，把数值粘性与方法误差分开（目前 R2 类结论绑定单一格式）。
\item 示踪点间距收敛扩展到 DNS 算例，多次实现（不同种子）给出误差分布。
\item 可压缩非线性精确解（等熵涡与声波相互作用、强激波熵产生）。
\item 与 LHD 报告一对一：按 \texttt{run\_rm.py} 重跑 RM；补 T4 同边界线积分、T7 动态 Laplacian；给爆炸波补初始密度；成对拟合改按 $K(t)$ 邻接挑边。
\item 共模误差的外部约束（独立压力探头、已知边界条件）。
\item 三维（四面体账本已在制造映射上验证）与真实 PTV 数据。
\end{enumerate}

\section{局限}
全部为二维；DNS 只有一档网格，误差多为单次实现；R2 参照与示踪点同源；压力需要参考常数；加速度类通道要求 $N_a\lesssim0.05$；第二定律只给上下界；粘度反演要求网络边界无涡量；三点映射耗散对噪声与曲率敏感；混合质量接管依赖 Pe 并受信息极限约束；"导热削弱压缩"依赖冷重气体构型。

\appendix
\section{代码与数据}
\begin{itemize}
\item \texttt{ltn\_dyn.py}：LTN 类（继承 \texttt{trajnet.TrajectoryNetwork}），全部通道的统一入口。
\item 掺混：\texttt{netmix.py}，\texttt{t1\_sod\_filter.py} 至 \texttt{t5c.py}，\texttt{rm\_*.py}。
\item 动力学通道：\texttt{hodge\_dyn.py}，\texttt{visc\_hodge.py}，\texttt{dissip\_tpm.py}，\texttt{entropy*.py}。
\item 声学：\texttt{ac\_run.py}，\texttt{ac\_dirac.py}；调和：\texttt{harm\_lo.py}；噪声：\texttt{gcv\_ell.py}，\texttt{cross\_dirichlet.py}。
\item 精确解：\texttt{suite.py}，\texttt{suite2.py}，\texttt{viz.py}；DNS 敏感性：\texttt{verify.py}，\texttt{tg\_verify.py}。
\item LHD 算例重算：\texttt{allcases.py}，\texttt{viz\_recalc.py}，\texttt{fig21.py}；基线：\texttt{compare.py}。
\end{itemize}
\end{document}
